% GENERATED FILE. Edit canonical lessons, then run npm run generate:books.
% canonical-source-hash: fnv1a64:48f43214924e738c
% canonical-lessons: ES-C56-cion

\chapter{-cion --- About Two Thousand Words, At Once}
\label{ch:friend-cion}
\hlchaptermodality{23}

\begin{chapteropening}
\textbf{By the end of this chapter:} \emph{I can turn an English -tion word into Spanish and know its gender without being told.}

Complete ES-C56-cion: convert an English -tion word, give the gender, and say why the correspondence is a rule.
\end{chapteropening}

\section[-ción]{-ción — about two thousand words, at once}

\label{lesson:ES-C56-cion}



\begin{quote}
You have learned roughly twenty Spanish words so far. This chapter is going to hand you about two thousand more, and it is not a trick.
\end{quote}

\begin{sounds}
\begin{itemize}
\raggedright
  \item \texttt{ci-th-s} — \emph{ci} is a soft sound: \emph{th} in most of Spain, \emph{s} across Latin America. \emph{Nación} is \emph{na-SYON} or \emph{na-THYON}, never \emph{na-KYON}.
\end{itemize}
\end{sounds}

\begin{grammarlens}[title={Grammar Lens: an ending that translates itself}]
\begin{quote}
\textbf{Spanish \emph{-ción} is English \emph{-tion}.}
\end{quote}

Not \textquotedblleft{}often.\textquotedblright{} Not \textquotedblleft{}sometimes.\textquotedblright{} As a rule:

\begin{quote}
\emph{nación} $\to$ nation · \emph{información} $\to$ information · \emph{educación} $\to$ education \emph{conversación} $\to$ conversation · \emph{nación} $\to$ nation · \emph{atención} $\to$ attention
\end{quote}

Read those without the English column and you can still read them.

\textbf{And every single one is feminine}, without exception — \emph{la nación}, \emph{la información}. One ending, one gender, no decision to make. That alone is worth the chapter: gender is the thing this book will otherwise ask you to learn word by word.

So: take almost any English word ending in \emph{-tion}, swap the ending for \emph{-ción}, and you have said something real. You will be right far more often than you are wrong, and when you are wrong you will still be understood.
\end{grammarlens}

\begin{cousinweb}
The reason this works is not luck, and it is worth knowing so you trust the rule rather than half-trusting it.

Latin had an ending \textbf{-tiōnem} that turned a verb into a noun for the act of doing it. Spanish \textbf{inherited} it directly, wearing it down to \emph{-ción}. English \textbf{borrowed} it — first from French after 1066, then straight from Latin for centuries afterwards — and kept it as \emph{-tion}.

Same ending, two routes, and the routes did not change it much. That is why the correspondence is a rule and not a resemblance: these are not similar words, they are \textbf{the same word}, arriving from two directions.
\end{cousinweb}

\subsection*{Your turn}
Say these aloud:

\begin{itemize}
\raggedright
  \item \textquotedblleft{}nación, información, educación\textquotedblright{} — hear the ending stay put
  \item an English \emph{-tion} word of your own, converted
\end{itemize}

\begin{tcolorbox}[breakable,colback=teal!4,colframe=teal!35!black,title={Before you move on}]
English \emph{-tion} becomes? (\emph{\textbf{-ción}}.) What gender? (\textbf{Feminine}, always.) And roughly how many words did that just give you? (About \textbf{two thousand}.) Next: another ending, and another twelve hundred.
\end{tcolorbox}
